\section{Propositions}
\label{section:proposition}

\subsection{Evaluation: Benchmarks, red teaming, and the skeptic}
\label{section:Evaluation}

\begin{summarybox}
We need to know whether the system actually works. This section argues for three complementary moves. First, a multidimensional evaluation that tells us \emph{where} the system fails rather than merely how often. Second, a red-teaming set of deceptive problems that expose specific failure modes such as retrieval mimicry, semantic collapse, and mode collapse. Third, a security and compliance evaluation to determine whether our product is ready for open release.
\end{summarybox}


\paragraph{Multidimensional, non-saturating evaluation.} The General Scales methodology \cite{Zhou2026} annotates benchmark instances along 18 capability demands with unbounded difficulty levels, giving an explanatory profile rather than a single score and remaining comparable across system generations as benchmarks would otherwise saturate. This might be an interesting direction for us. It tells us \emph{where} the system fails, not just how often. Paired with the QuantumQA generation pipeline \cite{2604.18176} (seeding$\rightarrow$generation to avoid mode collapse, hybrid solver+LLM-critic verification, task-adaptive answer formats), we could build a large annotated physics evaluation set spanning demand dimensions and difficulty levels.

\paragraph{A red-teaming set of deceptive problems.} The linear-drag droplet problem shared recently on Slack illustrated a particular trap: a non-trivial problem that looks like a standard textbook exercise. Frontier models exhibited characteristic failures on it: \emph{Retrieval Mimicry} (pattern-matching over reasoning, including ``interpolated'' derivations where steps don't actually follow from one another), and \emph{Semantic Collapse} (defaulting to the popularly discussed solution).

This is, admittedly, just one anecdote. What we need instead is a proper red-teaming set of deceptive problems. Arguably, they might be even more important than general benchmarks at this point. Currently, we do not even know, frankly, whether we compete on the raw power of \emph{solving} problems or on some other dimensions of user experience (more on this later). With a curated list of problems that target different error or failure modes (see e.g., the broad classification in the Psychopathia Machinalis \cite{Watson2025}\footnote{%
    The paper proposes a ``synthetic nosology'' for machine psychology, cataloging 32 maladaptive AI behaviors across seven axes and grouping failures into clinical-style syndromes; it offers a shared failure-mode vocabulary rather than evidence for any architecture choice.
    %Full citation: N.~Watson and A.~Hessami, \emph{Psychopathia Machinalis: A Nosological Framework for Understanding Pathologies in Advanced Artificial Intelligence}, Electronics 14, 3162 (2025). \url{https://doi.org/10.3390/electronics14163162}.
    },
and the recent diversity-collapse work \cite{2605.26492}), we could try to measure distinct error modes and subsequently design various interventions to improve in those dimensions.
By making our model more skeptical, more creative, or more ``scientific'' in other respects, we could already win the hearts of many users who do not appreciate the general "feel" that most major players offer (chat-heavy, too agreeable, etc.). Those targeted interventions -- to improve the product where the major competitors fall short, to optimize for a specific group of users -- while OpenAI and Anthropic must appeal to a broader audience -- might be our winning strategies.

\paragraph{Integrity and evaluation validity.} SciIntBench \cite{2605.29468} shows integrity alignment is strongly framing-sensitive, namely, models that refuse overt misconduct comply with covert framings, which matters for any user-facing research product. In addition, \cite{frasertaliente2026natural}\footnote{
    Anthropic's Natural Language Autoencoders (NLA) train a paired activation verbalizer and reconstructor to round-trip a residual-stream activation through a natural-language description, giving a readable interface to model internals; NLA-equipped auditing agents found a model organism's root misalignment, and the work reports models are evaluation-aware far more often (26\%) than their chain-of-thought verbalizes (1\%).
    %Full citation: Fraser-Taliente, Kantamneni, Ong et al., \emph{Natural Language Autoencoders Produce Unsupervised Explanations of LLM Activations}, Transformer Circuits Thread (2026). \url{https://transformer-circuits.pub/2026/nla/index.html}.
    }
shows models are evaluation-aware far more often than their chain-of-thought verbalizes (26\% vs 1\%), which can silently bias our own benchmarking. While the technique presented in the paper is too expensive for production monitoring, it can be used as a debugging instrument to examine unexpected agent behaviors.

This evaluation section has two further topics that deserve to be made explicit. The first is security. A research product that ingests literature, runs tools, and maintains long sessions is exposed to both direct and time-delayed \cite{2507.07341, 2510.01529}\footnote{
    The authors show that, under standard cryptographic assumptions (e.g. the existence of Time-Lock Puzzles, or symmetric keys planted in the model), there exist LLMs for which no efficient black-box filter can reliably separate harmful from benign prompts or outputs, provided the filter is computationally weaker than the model.
    %Full citation: Sarah Ball et al. \emph{On the Impossibility of Separating Intelligence from Judgment: The Computational Intractability of Filtering for AI Alignment}, arXiv:2507.07341 (2025). \url{https://arxiv.org/abs/2507.07341}.
    \\
    This is demonstrated as a "controlled-release prompting". The authors encode malicious prompts in one of two ways: timed-release (simple substitution ciphers that read as gibberish) or spaced-release (each character is replaced by verbose descriptive sentences that read as fluent natural language). The lightweight guard lacks the inference budget or context window to decode the prompt and passes it through; the main LLM has the resources to follow the decoding instructions and execute the hidden jailbreak.
    The authors reported "near-perfect" success rates on Gemini 2.5 Flash, DeepSeek DeepThink, Grok 3, and Mistral/Magistral. OpenAI, Anthropic, and Meta AI demonstrated robust resistance, attributed to output filtering and training-based alignment that goes beyond input monitoring.
    %Full citation: Jaiden Fairoze et al., \emph{Bypassing Prompt Guards in Production with Controlled-Release Prompting.} arXiv:2510.01529 (2025). \url{https://arxiv.org/abs/2510.01529}.
    }
prompt injection, so we need to monitor both the input \emph{and} the output distribution (and possibly make the prompt guards powered by an LLM of equal or \emph{larger} size than the underlying reasoning model).

Another aspect is compliance, which is elaborated further in Section \ref{section:compliance}. In short, to preserve the limited-risk classification under the EU AI Act, we should develop a harness that prevents users from working with clinical data and implement CBRN-relevant mitigation if our system scores highly on those benchmarks. To avoid systemic design flaws that could stall our market launch and cause costly architectural rework, we should address these risks proactively during the design phase, not as afterthoughts.

\subsection{Data, signal, and user sessions}
\label{section:data_signal_sessions}

\begin{summarybox}
Our data and the shape of our users' sessions could be how we differentiate ourselves from other players that offer more general, consumer-focused chat interfaces. Three threads matter here. First, training signal: the most valuable signal may come from verified edits, interventions, corrections, and re-runs rather than generic human feedback. Second, memory: scientific sessions span days and involve large bodies of literature, where naive retrieval and context-management techniques may be insufficient. Third, the platform: if designed carefully, cached summaries, cross-session learning, and network effects could become a compounding advantage.
\end{summarybox}


\paragraph{Training signal and data.}
The QuantumQA \cite{2604.18176} result that matters most strategically is this: the majority of model errors were logical/physical, not calculational, and the solver-grounded RL reduced only the logic/physics errors, while calculation errors remained flat. 
This gives us an opportunity.
Our symbolic-representation and symbolic-manipulation workstreams attack the bucket that the QuantumQA's authors couldn't move.
Another finding from QuantumQA deserves a mention: the negative correlation between token length and solution quality, with standard RLHF rewarding verbose plausibility. This is a warning that even a large user base producing human feedback might be less useful for training than physics-grounded verification signals. This should be verified and discussed with Product. It is likely that the interactions worth capturing are edits, interventions, corrections, and re-runs validated against verifiers, not thumbs-up ratings.

\paragraph{Context and memory in long sessions.}
Our users' sessions are unlike consumer chat: if we want to tap into a scientific process, they should likely span many days, with semi-persistent context and reasoning over a large body of literature. This creates challenges that might not be present in customer-focused products from OpenAI and Anthropic. This might be our chance to address them.

Retrieval-based systems lose precision as chunk counts grow, while models holding full documents in context achieve near-perfect retrieval \cite{2508.21038}\footnote{
    This work proves that single-vector embedding models face a dimension-bounded ceiling (via sign-rank) on which top-$k$ document combinations they can return, while a long-context reranker fed all documents at once achieves near-perfect retrieval only on small configurations.
    %Full citation: O.~Weller, M.~Boratko, I.~Naim, J.~Lee, \emph{On the Theoretical Limitations of Embedding-Based Retrieval}, arXiv:2508.21038 (2025). \url{https://arxiv.org/abs/2508.21038}.
    }.
Detailed paper summarization (e.g., ARA creation, \cite{2604.24658}) may cost up to a dollar per paper (when used with a capable model). The costs of summarizing all of the available data could be prohibitive. 
An alternative could be a fast query-conditioned summarization. As an example, DiffusionGemma is a new $>$1000 tok/s diffusion model \cite {google2026diffusiongemma}. Accepting lower quality, we could compress each paper 20$\times$ faster and 1-2 orders of magnitude cheaper. This allows us to produce such a summary on demand \emph{for the specific question at hand}, enabling reasoning over dozens or hundreds of papers at once (e.g., with the option to retrieve the original text to verify the claims).

\paragraph{Cross-instance learning.}
We could also treat our platform as a platform and leverage cross-instance learning. Every computed summary can be cached, and every discovered error in a paper or cross-paper link can be stored and reused. We might also consider cross-session learning (of course, with some privacy mitigation strategies in place). This is a genuine compounding advantage: more users means faster retrieval and more accurate reasoning over the literature, something a general-purpose provider cannot match in our domain. One important caveat: privacy boundaries between local and shared knowledge must be designed from the start. We cannot say ``this is closed beta, let's worry about it later'', reworking the solution architecture later would be very hard and costly.

\paragraph{Long-term stability.}
For session stability, we might leverage a wiki-like local memory with branching: the system explores alternative branches, writes a post-mortem when a branch dead-ends, commits the lesson to the wiki, and backtracks.\footnote{This plays well with the concept of the dynamical research objects (DRO). Not fully, but it could be a good starting point.} A global wiki would accumulate domain-general lessons across sessions and users. While unbounded context growth, we might introduce retiring aspects (forgetting) as described in AutoResearchClaw \cite{2605.20025}. It remains an unexplored topic, and we could make original contributions here.\footnote{
    There is one recent paper that already discusses this topic, by introducing AgingBench \cite{2605.26302}, a longitudinal benchmark that measures how deployed agents degrade over multiple sessions. They organize aging into four mechanisms: compression (loss of detail during write-time summarization), interference (confusion from accumulated similar memories at retrieval), revision (failure to update derived or changed state), and maintenance (regression after lifecycle events such as memory flushing/recompaction).
    %Full citation: Jianing Zhu et al., \emph{Your Agents Are Aging Too: Agent Lifespan Engineering for Deployed Systems}, arXiv:2605.26302 (2026). \url{https://arxiv.org/abs/2605.26302v1}.
}

\paragraph{Personalization and network effects.}
A different question is how much this should result in ``automatic customizability'' to particular users. It is tempting, but it has privacy implications and could push us into the high-risk category in Europe (see Section \ref{section:compliance}). There is also the risk of putting scientists into information bubbles (presenting only topics that reinforce their biases), which we want to avoid.

Building our infrastructure, we will be able to eventually tap into the network effects: e.g., pairing researchers with shared interests, as done in SciMuse \cite{2405.17044}.\footnote{
    SciMuse generates personalized interdisciplinary research ideas by building a 58-million-paper knowledge graph and prompting GPT-4 with concept pairs, and is evaluated by 110 Max Planck group leaders who rate 4,451 suggestions on a 1–5 interest scale. The headline result is split: the knowledge graph gave no measurable lift to idea generation over a pure-LLM baseline (a null with no reported test statistic), but graph features powered the best selector in the paper, beating zero-shot GPT-4o on top-N precision ($\approx$66\% vs 45\% at top-3).
    %Full citation: X.~Gu and M.~Krenn, \emph{Interesting Scientific Idea Generation using Knowledge Graphs and LLMs: Evaluations with 100 Research Group Leaders}, arXiv:2405.17044 (2024).
    }
Another interesting mode would be to allow group collaboration on a problem across different branches of the exploration tree. These are later-stage opportunities, contingent on having a user network at all, but the initial design should not preclude them.

\subsection{Economic and regulatory constraints}
\label{section:compliance}

\begin{summarybox}
A strong system can still lose on economics or stumble on regulation, so we treat both as design inputs. The economic trap is that the person using the product and the person paying for it may be different personas. Individual researchers may compare us against frontier-model access that feels free, while institutions will care about total cost, privacy, control, and workflow fit. On the regulatory side, the for-profit pivot removes the EU AI Act research exemptions; we will likely fall into the limited-risk class, but clinical-data use and CBRN-relevant capability could create higher-risk obligations. These are design constraints, not afterthoughts.
\end{summarybox}


\paragraph{Economics of the comparison.}
Targeting scientists, we might be at a structural disadvantage. Professors are constrained in how they can use their research budgets. Some universities currently give faculty effectively unlimited access to AI. Researchers from those institutions might perceive frontier-model queries as having a cost near zero. The harsh truth is that even a 100$\times$ cost reduction cannot compete with something that is given for free. Two consequences. First, head-to-head ``us vs.\ frontier model'' comparisons made by academics will likely be one-sided; they never see the cost axis on which we can win -- however, they will care about the quality and user experience. Second, the economics change entirely if the customer is the university rather than the researcher: institutions do care about total spend, and our domain-specific optimizations (caching, summaries, smaller models, persistent cross-user knowledge) translate directly into a cost story at that level.

If we serve our own models, we might also try to initially compete on trust and fit. A model running locally (or on designated on-prem infrastructure on the client's side) addresses both the privacy and control concerns some institutions (like national labs) might have.
This and a better UX, one genuinely suited to scientific workflows, e.g., with the ability for researchers to collaborate within the same session, branch the project, version control ideas, etc., might be the biggest differentiator for us.

\paragraph{Regulatory constraints on the design.}
With the pivot to for-profit, the EU AI Act research exemptions (development-phase and scientific use) no longer protect us. The likely classification is limited-risk (transparency obligations), which is manageable with documentation, guardrails, and monitoring artifacts. To maintain this category, we might need to develop a harness that prevents users from using our product with clinical data (so we do not drift into the high-risk category that might require audits, etc.). In addition, if the system scores highly on CBRN-relevant benchmarks, as is plausible for a capable science system, we might need to introduce mitigations to reduce our capability in selected topics. Filters, guardrails, and selective unlearning are all tricky, and unlearning in particular is often reversible. This is a real (and hard) problem, and we need to discuss how we will approach it for real.

\subsection{Self-improvement: The theoretical frame}
\label{section:self_improvement}

\begin{summarybox}
Here, we explore different starting points for designing a self-improving system. The physics-verification bet gives us hope that such a direction might be sustainable: self-improvement loops are fragile when they feed only on their own synthetic outputs, but they may remain stable when anchored to external signals such as solvers, proof checkers, code execution, or experimental data.
\end{summarybox}


The Darwin--Gödel machine line of work \cite{2505.22954}
might be a good starting point if we decide to build self-improving loops. Gödel's incompleteness implies a system cannot, in general, prove which self-modifications are beneficial; the authors proposed a Darwinian-like selection under environmental pressure to continuously improve (evolve) an agentic system  -- there is no guarantee of optimality, only an effective stochastic search. Whether selection-type pressure is \emph{efficient} for AI self-improvement remains an open question, with arguments for and against in the literature. ERA \cite{Aygn2026} shows a similar approach in practice, blending genetic programming and tree-search ideas.
The theoretical limits matter too \cite{2601.05280} formalizes how a closed self-improvement loop that refits its own unfiltered samples collapses as the share of fresh external signal vanishes. However, importantly, the authors also say:
\begin{quotation}
    \noindent ``Empirical studies indicate that collapse depends critically on the proportion of real versus synthetic
    data, with degradation emerging when synthetic data dominate and stability maintained when sufficient external signal persists.''
\end{quotation}
Thus, by providing external grounding, whether in the form of verifiable feedback (physics with a solver, math with proof checkers, code with execution) or comparison with experimental data, the system can sustain improvements and avoid degenerative dynamics. That is encouraging for our physics-verification bet.

\clearpage
\subsection{Advanced architectural design}
\label{section:advanced_architecture}

\begin{summarybox}
In this section, we argue that good architecture and good UX, even if perfectly observable, may still be hard to copy. The differentiator is not only the visible architecture, but the ``glue'' between components: timing, interaction design, perceived speed, memory, recovery, and the way the system adapts to the user's actual intent. For our users, accuracy matters, but so do speed, cost, usability, transparency, and workflow fit.
\end{summarybox}

One perspective is that architecture is easy to copy. That may hold for something like AxProverBase \cite{2602.24273}: they proposed a simple agentic flow that is easy to replicate, including by us. But, as discussed with Melissa, this notion is sometimes an illusion. The proper system design -- the whole user experience -- is not easy to replicate. Think of Google Search: the page is minimalistic, yet the experience is unmatched by similar products. Its strength is speed, precision, and minimalism -- people see what matters to them, no more and no less.

In our current system, users have to wait a long time. Instagram's trick was to start preloading images (or videos) while people were busy editing the photo --  cropping, picking a filter, typing a caption. The perceived performance was far higher than what was technologically possible: an illusion of near-instant upload.

\paragraph{Proposition 1.} While we present the results of an ordinary search, we run a detailed agentic search in the background. The user can review the first sources while we work on the more time-consuming task of bringing in more relevant papers. The early signal -- what the user selects, what they reject -- can further refine the search criteria, since there is often a gap between what a user asks for and what they actually want.

Another pain point is the static definition of our phases. It creates a difficult entry point for a user accustomed to the near-instant question-and-response loops of Anthropic, OpenAI, and Google. A more flexible definition of phases could be more appropriate. But the rigid structure has advantages too: we can optimize tool usage, tailor the experience, and go much deeper precisely because we know what the user wants. In addition, being somewhat rigid is a way to be \emph{different} from other products, which may appeal to scientists who have found that research conducted with monolithic LLMs can be shallow or superficial. A solution to this tension would be a hybrid approach, discussed in the next propositions.

\paragraph{Proposition 2.} Use an architect agent to define the flow, and a classical executor to build the pipeline (a DAG, with all loops expanded flat). Elements are either selected from a repository of templates or constructed on the fly. An operational wiki is used throughout, and a post-mortem mechanism suggests recovery when something fails.

\paragraph{Proposition 3.} To further leverage cross-session learning, a hierarchical wiki is used: one local and operational, one global and cross-run.

\paragraph{Proposition 4.} Context can overflow. Thus, treat the wiki as a set of statements drawn from a distribution with a half-life of about 30 days (or another value). This allows exploration and recovery from bad memories, via either periodic pruning of old statements or a refresher mechanism, combined with the ability to rewrite or correct a memory.